% Part: methods
% Chapter: proofs
% Section: cant-do-it

\documentclass[../../../include/open-logic-section]{subfiles}

\begin{document}

\olfileid{mth}{prf}{cnt}

\olsection{আমি পারছি না!{}}

আমাদের সবারই কখনো মনে হয়, ছেড়ে দিই। কিন্তু তুমি
\emph{পারবে}। তোমার শিক্ষক, সহশিক্ষক এবং সহপাঠীরা
সাহায্য করতে পারেন। তাঁদের কাছে সাহায্য চাও!{}
হাল ছেড়ে দেওয়ার পরিস্থিতি এড়াতে, আর তেমন মনে
হলে কী করতে পারো, তার কয়েকটি পরামর্শ নিচে।

সমস্যা সফলভাবে সমাধান করতে এগুলি করো:
\begin{enumerate}
\item \emph{যত আগে সম্ভব শুরু করো।} পুরো সেমিস্টারে
  নানা কাজ থাকে; অনেকেরই কাজ ফেলে রাখার অভ্যাস
  আছে। তাই বাড়ির কাজ আগে শুরু করা খুব উপকারী।
  আটকে গেলে সমাধান খোঁজার সময় পাবে—কান্নায়
  ভেঙে পড়া ছাড়া আরও পথ থাকবে।
\item \emph{সহপাঠীদের সঙ্গে কথা বলো।} তুমি একা নও।
  অন্যরাও আটকে যেতে পারে, হয়তো ভিন্ন জায়গায়।
  সহপাঠীর সঙ্গে আলোচনা করলে সমস্যাটি অন্যভাবে
  দেখতে পারো; সেখান থেকে সমাধানের পথ খুলতে পারে।
  অবশ্য তাদের সমাধান নকল কোরো না: ইঙ্গিত চাও,
  কোথায় আটকে গেছ বোঝাও এবং পরের ধাপটি জানতে
  চাও। নিজে বুঝলে অন্যকেও সাহায্য করো। অন্যকে
  বোঝাতে গেলে নিজের বোঝাও পোক্ত হয়।
\item \emph{সাহায্য চাও।} তোমার জন্য নানা সহায়তা
  আছে। শিক্ষক ও সহশিক্ষক তোমাকে সাহায্য করতে
  আছেন এবং তুমি সফল হও, তা \emph{চান}। সমস্যাটি
  কোথায় আটকে যাচ্ছে, তা বুঝে এগোতে তাঁরা
  সাহায্য করতে পারবেন।
\item \emph{একটু বিরতি নাও।} আটকে যাওয়ার কারণ
  \emph{হতে পারে}, অনেকক্ষণ ধরে একই সমস্যা
  দেখছ। অল্প বিরতি নাও, চা খাও, বা কিছুক্ষণ
  অন্য সমস্যা নিয়ে কাজ করো। তারপর নতুন মনে
  ফিরে এসো। দরকার হলে এক রাত ঘুমিয়ে নাও।
\end{enumerate}

খেয়াল করো, এই কৌশলগুলি কাজে লাগাতে অনেক আগে
থেকে প্রমাণটি নিয়ে কাজ শুরু করতে হয়। জমা দেওয়ার
আগের রাত দুটোয় শুরু করলে এগুলি তত কাজে নাও
আসতে পারে।

শুনতে হতাশাজনক লাগতে পারে, কিন্তু প্রমাণের সমস্যা
সমাধান করা একটি ফলপ্রসূ চ্যালেঞ্জ। কেউ কেউ
এটিকেই পেশা করেন—নিশ্চয়ই এতে আনন্দেরও কিছু
আছে। প্রায় সব কিছুর মতো সমস্যা সমাধান ও প্রমাণ
করার জন্যও অনুশীলন দরকার। কোনো সহপাঠীর কাছে
সহজ মনে হলে, সম্ভবত সে ইতিমধ্যে অনেক অনুশীলন
করেছে। খুব তাড়াতাড়ি হাল ছেড়ো না।

নির্দিষ্ট কোনো সমস্যায় সময় বা ধৈর্য ফুরিয়ে গেলে
অসুবিধা নেই। তার মানে তুমি বোকা নও, বা কোনো
দিনই বুঝবে না—এমন নয়। শিক্ষক বা সহপাঠীর কাছে
জেনে নাও, কাজটি কীভাবে হয়। নিজের ভুল বা আটকে
যাওয়ার জায়গাটি চিহ্নিত করো, যাতে একই ধরনের
সমস্যায় আবার না আটকে যাও। তারপর সমাধান না
দেখে নিজে চেষ্টা করো। পরের বার আরও আগে
শুরু করো এবং আগে সাহায্য চাও।

\end{document}
